\documentclass[11pt, a4paper]{article}
\usepackage[a4paper, margin=2.5cm]{geometry}
\usepackage[utf8]{inputenc}
\usepackage[T1]{fontenc}
\usepackage[spanish,es-nodecimaldot]{babel}
\usepackage{enumitem}
\setlist[itemize]{label=-}
\usepackage{amsmath}
\usepackage{amssymb}
\usepackage{booktabs}
\usepackage{xcolor}
\usepackage{listings}
\usepackage{hyperref}

\makeatletter
\setlength{\@fptop}{0pt}
\setlength{\@fpbot}{0pt plus 1fil}
\makeatother

\lstset{
    language=Matlab,
    backgroundcolor=\color{gray!6},
    basicstyle=\ttfamily\footnotesize,
    keywordstyle=\color{blue}\bfseries,
    commentstyle=\color{green!50!black}\itshape,
    stringstyle=\color{purple},
    numbers=left,
    numberstyle=\tiny\color{gray},
    stepnumber=1,
    numbersep=8pt,
    frame=single,
    rulecolor=\color{gray!30},
    tabsize=4,
    showspaces=false,
    showstringspaces=false,
    breaklines=true,
    breakatwhitespace=true,
    literate={á}{{\'a}}1 {é}{{\'e}}1 {í}{{\'i}}1 {ó}{{\'o}}1 {ú}{{\'u}}1 {ñ}{{\~n}}1 {Á}{{\'A}}1 {É}{{\'E}}1 {Í}{{\'I}}1 {Ó}{{\'O}}1 {Ú}{{\'U}}1 {Ñ}{{\~N}}1
}

\title{\textbf{Relación entre Erlang B y Erlang C}}
\author{David Stiven Muñoz Amaya - 20202020071}
\date{}

\begin{document}

\maketitle

\section{Interpretación conceptual}

Los sistemas Erlang B ($M/M/k/k$, con pérdida) y Erlang C ($M/M/k/\infty$, con espera) se comportan de forma idéntica mientras existan servidores libres ($0 \le n < k$), ya que la tasa de llegada es $\lambda$ y la de servicio es $n\mu$. Por ello, ambos comparten los mismos términos proporcionales $\frac{A^n}{n!}$, donde $A = \frac{\lambda}{\mu}$ es el tráfico ofrecido.

La diferencia ocurre al alcanzar el límite de $k$ servidores ocupados:
\begin{itemize}
    \item \textbf{En Erlang B:} El sistema se trunca en $n = k$. Todo cliente que llega cuando $N = k$ es bloqueado y se pierde.
    \item \textbf{En Erlang C:} Si llegan más clientes ($n \ge k$), forman una fila infinita atendida a tasa constante $k\mu$. Al sumar las probabilidades de todos los estados de espera ($n \ge k$) mediante una serie geométrica (con $A < k$), aparece el factor de amplificación de la cola $\left(\frac{k}{k - A}\right)$ acompañando al término $\frac{A^k}{k!}$.
\end{itemize}

\section{Demostración matemática}

Para el sistema Erlang B, la probabilidad de bloqueo $E_{1,k}$ está dada por:
\[
E_{1,k} = \frac{\frac{A^k}{k!}}{\sum\limits_{n=0}^{k} \frac{A^n}{n!}}
\]

Separando el último término ($n = k$) de la sumatoria del denominador e invirtiendo la fracción:
\[
\frac{1}{E_{1,k}} = \frac{\sum\limits_{n=0}^{k-1} \frac{A^n}{n!} + \frac{A^k}{k!}}{\frac{A^k}{k!}} = \frac{\sum\limits_{n=0}^{k-1} \frac{A^n}{n!}}{\frac{A^k}{k!}} + 1
\]

Despejando el término de la sumatoria:
\begin{equation}
\frac{\sum\limits_{n=0}^{k-1} \frac{A^n}{n!}}{\frac{A^k}{k!}} = \frac{1}{E_{1,k}} - 1 = \frac{1 - E_{1,k}}{E_{1,k}}
\end{equation}

Por otro lado, para el sistema Erlang C ($M/M/k/\infty$), la probabilidad de espera se estructura como:
\[
E_{2,k} = P[T_w > 0] = \frac{A^k}{k!} \left( \frac{k}{k - A} \right) P_0
\]
donde la probabilidad de sistema vacío ($P_0$) viene dada por:
\[
P_0 = \left[ \sum\limits_{n=0}^{k-1} \frac{A^n}{n!} + \frac{A^k}{k!} \left( \frac{k}{k - A} \right) \right]^{-1}
\]
Al sustituir $P_0$ en la expresión de $E_{2,k}$, se obtiene la ecuación completa de Erlang C:
\[
E_{2,k} = \frac{\frac{A^k}{k!} \left(\frac{k}{k - A}\right)}{\sum\limits_{n=0}^{k-1} \frac{A^n}{n!} + \frac{A^k}{k!} \left(\frac{k}{k - A}\right)}
\]


Dividiendo el numerador y el denominador de la expresión de $E_{2,k}$ entre $\frac{A^k}{k!}$:
\[
E_{2,k} = \frac{\frac{k}{k - A}}{\frac{\sum\limits_{n=0}^{k-1} \frac{A^n}{n!}}{\frac{A^k}{k!}} + \frac{k}{k - A}}
\]

Sustituyendo la ecuación (1) en el denominador:
\[
E_{2,k} = \frac{\frac{k}{k - A}}{\left(\frac{1 - E_{1,k}}{E_{1,k}}\right) + \frac{k}{k - A}}
\]

Multiplicando el numerador y el denominador por $E_{1,k}(k - A)$ para simplificar las fracciones:
\[
E_{2,k} = \frac{k \cdot E_{1,k}}{(1 - E_{1,k})(k - A) + k \cdot E_{1,k}}
\]

Expandiendo el producto en el denominador y cancelando los términos $-k \cdot E_{1,k}$ y $+k \cdot E_{1,k}$:
\[
\text{Denominador} = k - A - k \cdot E_{1,k} + A \cdot E_{1,k} + k \cdot E_{1,k} = k - A + A \cdot E_{1,k}
\]

Factorizando $A$ en el denominador, se obtiene finalmente la ecuación de relación directa:
\[
\mathbf{E_{2,k} = \frac{k \cdot E_{1,k}}{k - A + A \cdot E_{1,k}} = \frac{k \cdot E_{1,k}}{k - A(1 - E_{1,k})}}
\]

Donde:
\begin{itemize}
    \item $E_{1,k}$: Probabilidad de bloqueo en Erlang B.
    \item $k$: Número de servidores ($k \in \mathbb{Z}^+$).
    \item $A$: Intensidad de tráfico ofrecido $A = \frac{\lambda}{\mu}$ (bajo condición de estabilidad $A < k$).
\end{itemize}

\section{Implementación en MATLAB}

\subsection{Función base de Erlang B (\texttt{pb.m})}

\begin{lstlisting}
function p = pb(A, k)
    p = A ./ (1 + A);
    for n = 2:k
        p = A .* p ./ (n + A .* p);
    end
end
\end{lstlisting}

\subsection{Función de Erlang C (\texttt{erlangC.m})}

\begin{lstlisting}
function E2k = erlangC(A, k)
    if any(A <= 0 | A >= k)
        error('A debe ser positivo y menor que k (0 < A < k).') 
    end
    if any(k <= 0 | mod(k, 1) ~= 0)
        error('k debe ser un numero entero positivo.')
    end

    E1k = pb(A, k);

    E2k = (k .* E1k) ./ (k - (A .* (1 - E1k)));
end
\end{lstlisting}

\section{Verificaciones}

\begin{lstlisting}
A_base = [0.2, 0.5, 0.8];
k_base = 1;

    for i = 1:length(A_base)
        P_esp = erlangC(A_base(i), k_base);
        fprintf('A = %4.1f | k = %3d | P(espera) = %8.4f\n', ...
                A_base(i), k_base, P_esp);
    end


pruebas = [
    2.0,   3;
    2.0,   5;
    2.0,  10;
    4.0,   5;
    4.0,  10;
    4.0,  20;
    4.0,  40;
    4.0, 200
];

    for i = 1:size(pruebas, 1)
        A_val = pruebas(i, 1);
        k_val = pruebas(i, 2);
        P_esp = erlangC(A_val, k_val);
        fprintf('A = %4.1f | k = %3d (Cumple A < k) -> P(espera) = %.4e\n', ...
                A_val, k_val, P_esp);
    end
\end{lstlisting}

\clearpage
\begin{table}[t]
\centering
\caption{Resultados de la verificación de la probabilidad de espera para diferentes valores de $A$ y $k$.}
\label{tab:resultados_verificacion}
\vspace{0.3em}
\begin{tabular}{ccccc}
\toprule
\textbf{$A$} & \textbf{$k$} & \textbf{Condición $A < k$} & \textbf{$P[T_w > 0]$ (Científica)} & \textbf{$P[T_w > 0]$ (\%)} \\
\midrule
0.2 & 1   & Sí ($0.2 < 1$)   & $2.0000 \times 10^{-1}$   & $20.00\%$ \\
0.5 & 1   & Sí ($0.5 < 1$)   & $5.0000 \times 10^{-1}$   & $50.00\%$ \\
0.8 & 1   & Sí ($0.8 < 1$)   & $8.0000 \times 10^{-1}$   & $80.00\%$ \\
\midrule
2.0 & 3   & Sí ($2 < 3$)    & $4.4444 \times 10^{-1}$   & $44.44\%$ \\
2.0 & 5   & Sí ($2 < 5$)    & $5.9701 \times 10^{-2}$   & $5.97\%$ \\
2.0 & 10  & Sí ($2 < 10$)   & $4.7737 \times 10^{-5}$   & $0.0048\%$ \\
\midrule
4.0 & 5   & Sí ($4 < 5$)    & $5.5411 \times 10^{-1}$   & $55.41\%$ \\
4.0 & 10  & Sí ($4 < 10$)   & $8.8147 \times 10^{-3}$   & $0.8815\%$ \\
4.0 & 20  & Sí ($4 < 20$)   & $1.0347 \times 10^{-8}$   & $\approx 0\%$ \\
4.0 & 40  & Sí ($4 < 40$)   & $3.0153 \times 10^{-26}$  & $\approx 0\%$ \\
4.0 & 200 & Sí ($4 < 200$)  & $6.1194 \times 10^{-257}$ & $\approx 0\%$ \\
\bottomrule
\end{tabular}
\end{table}

\end{document}